\documentclass[11pt]{article}
\usepackage[ngerman]{babel}
\newcommand{\imgdir}{images/de}
% Shared layout for the Illuminati manual. Each manual-<lang>.tex loads this
% after \documentclass and sets \manuallang (babel name) and \imgdir.
% Build with XeLaTeX (see build.ps1); the fonts are the app's own.

\usepackage{fontspec}
\setmainfont{Literata}[
  Path=../../assets/fonts/literata/, Extension=.ttf,
  UprightFont=*-Regular, BoldFont=*-Bold, ItalicFont=*-Italic, BoldItalicFont=*-BoldItalic]
\setsansfont{AtkinsonHyperlegible}[
  Path=../../assets/fonts/atkinson/, Extension=.ttf,
  UprightFont=*-Regular, BoldFont=*-Bold, ItalicFont=*-Italic, BoldItalicFont=*-BoldItalic]
\newfontfamily\dyslexicfont{OpenDyslexic}[
  Path=../../assets/fonts/opendyslexic/, Extension=.otf,
  UprightFont=*-Regular, BoldFont=*-Bold, ItalicFont=*-Italic, BoldItalicFont=*-Bold-Italic]
\newfontfamily\garamondfont{EBGaramond}[
  Path=../../assets/fonts/ebgaramond/, Extension=.ttf,
  UprightFont=*-Regular, BoldFont=*-Bold, ItalicFont=*-Italic, BoldItalicFont=*-BoldItalic]

\usepackage[a4paper, margin=2.4cm, headheight=14pt]{geometry}
\usepackage{graphicx}
\usepackage[table]{xcolor}
\usepackage{booktabs, tabularx, array}
\usepackage{enumitem}
\usepackage{fancyhdr}
\usepackage{titlesec}
\usepackage[most]{tcolorbox}
\usepackage[hidelinks]{hyperref}

% Colours from the app's logo.
\definecolor{violet}{HTML}{34206A}
\definecolor{violetlight}{HTML}{EDE8F7}
\definecolor{gold}{HTML}{C9A227}
\definecolor{goldlight}{HTML}{FBF4DF}

\setlength{\parskip}{0.55em}
\setlength{\parindent}{0pt}
\setlist{itemsep=0.25em, topsep=0.3em}
\linespread{1.08}

\titleformat{\section}{\sffamily\Large\bfseries\color{violet}}{\thesection}{0.8em}{}
\titleformat{\subsection}{\sffamily\large\bfseries\color{violet}}{\thesubsection}{0.7em}{}
\titlespacing*{\section}{0pt}{1.6em}{0.6em}
\titlespacing*{\subsection}{0pt}{1.1em}{0.35em}
% Don't leave a heading alone at the bottom of a page.
\usepackage{needspace, etoolbox}
\pretocmd{\section}{\needspace{8\baselineskip}}{}{}
\pretocmd{\subsection}{\needspace{5\baselineskip}}{}{}

\pagestyle{fancy}
\fancyhf{}
\fancyhead[L]{\sffamily\small\color{violet} Illuminati}
\fancyhead[R]{\sffamily\small\color{gray}\nouppercase{\rightmark}}
\fancyfoot[C]{\sffamily\small\thepage}
\renewcommand{\headrulewidth}{0.3pt}
% \rightmark is the first section mark on the page.
\renewcommand{\sectionmark}[1]{\markboth{#1}{#1}}
\renewcommand{\subsectionmark}[1]{}

% Keep catalogue names whole in every language.
\babelhyphenation{Wiki-source Guten-berg}

% A key or key combination, e.g. \kbd{Ctrl+F}.
\newcommand{\kbd}[1]{\tikz[baseline=(k.base)]\node[draw=gray!60, fill=gray!8,
  rounded corners=2pt, inner sep=2pt, font=\sffamily\small](k){#1};}
\usepackage{tikz}

% Arrow keys (the sans font has no arrow characters).
\newcommand{\kright}{\kbd{\ensuremath{\rightarrow}}}
\newcommand{\kleft}{\kbd{\ensuremath{\leftarrow}}}
\newcommand{\kup}{\kbd{\ensuremath{\uparrow}}}
\newcommand{\kdown}{\kbd{\ensuremath{\downarrow}}}

% A button, menu item or label as it appears in the app.
\newcommand{\ui}[1]{\textsf{\textbf{#1}}}

% Screenshot with caption; skipped if the image file isn't there.
\newcommand{\shot}[2]{%
  \IfFileExists{\imgdir/#1.jpg}{%
    \begin{figure}[htbp]\centering
      \fcolorbox{gray!40}{white}{\includegraphics[width=0.96\textwidth]{\imgdir/#1.jpg}}
      \caption{#2}
    \end{figure}}{}}

\newtcolorbox{tip}[1]{enhanced, colback=goldlight, colframe=gold,
  boxrule=0.6pt, arc=3pt, left=8pt, right=8pt, top=5pt, bottom=5pt,
  fonttitle=\sffamily\bfseries, coltitle=black, colbacktitle=gold!35,
  title={#1}}
\newtcolorbox{note}[1]{enhanced, colback=violetlight, colframe=violet,
  boxrule=0.6pt, arc=3pt, left=8pt, right=8pt, top=5pt, bottom=5pt,
  fonttitle=\sffamily\bfseries, coltitle=white, colbacktitle=violet,
  title={#1}}

% Two-column tables that fill the text width.
\newenvironment{keytable}[2]{%
  \renewcommand{\arraystretch}{1.3}%
  \rowcolors{2}{white}{violetlight}%
  \tabularx{\textwidth}{@{}>{\raggedright\arraybackslash}p{0.33\textwidth}X@{}}
  \toprule \textsf{\textbf{#1}} & \textsf{\textbf{#2}} \\ \midrule}{%
  \bottomrule\endtabularx}

% Title page: \maketitlepage{subtitle}{version line}{blurb}
\newcommand{\maketitlepage}[3]{%
  \begin{titlepage}
    \centering
    \vspace*{3cm}
    \includegraphics[width=4.2cm]{../../assets/icon/app_icon.png}\par\vspace{1.2cm}
    {\sffamily\fontsize{40}{46}\selectfont\bfseries\color{violet} Illuminati\par}
    \vspace{0.5cm}
    {\sffamily\LARGE #1\par}
    \vspace{1.5cm}
    {\large\itshape #3\par}
    \vfill
    {\sffamily\color{gray} #2\par}
  \end{titlepage}}

\hypersetup{pdftitle={Illuminati — Benutzerhandbuch}, pdfauthor={Illuminati}, pdflang={de}}

\begin{document}

\maketitlepage{Benutzerhandbuch}{Version 1.0 · Windows, macOS, Linux und Android}%
  {Finde, lade und lies Bücher, die nie urheberrechtlich geschützt waren\\oder deren Schutz abgelaufen ist.}

{\setcounter{tocdepth}{1}\tableofcontents}
\newpage

% ---------------------------------------------------------------------------
\section{Willkommen}

Illuminati ist eine Lese-App für Windows, macOS, Linux und Android. Sie funktioniert wie
E-Book-Apps, die du vielleicht schon kennst, mit einem Unterschied: Alle Bücher darin sind
kostenlos und legal herunterzuladen, weil ihr Urheberrecht abgelaufen ist oder nie bestand.

\begin{itemize}
  \item Über 75.000 Bücher von \textbf{Project Gutenberg}, dazu \textbf{Standard Ebooks},
        \textbf{Wikisource} und das \textbf{Internet Archive}.
  \item Ein angenehmer Reader mit Markierungen, Notizen, Lesezeichen, Suche, Wörterbuch und
        Übersetzung.
  \item Deine eigene Bibliothek mit Regalen und Leselisten.
  \item Anzeige- und Barrierefreiheitseinstellungen, darunter eine Schrift für Menschen mit
        Legasthenie.
  \item Kein Konto, kein Abo, keine Werbung. Alles bleibt auf deinem Gerät.
\end{itemize}

Die Oberfläche gibt es auf Deutsch, Englisch, Spanisch, Französisch, Italienisch und
Portugiesisch.

% ---------------------------------------------------------------------------
\section{Erste Schritte}

\subsection{Installation}
\begin{itemize}
  \item \textbf{Windows:} Kopiere den ganzen Ordner \texttt{Illuminati} auf deinen Computer und
        starte \texttt{illuminati.exe}. Lass die anderen Dateien daneben liegen; das Programm
        braucht sie.
  \item \textbf{Android:} Installiere die \texttt{.apk}-Datei. Eventuell musst du in den
        Einstellungen des Telefons die Installation aus dieser Quelle erlauben.
  \item \textbf{macOS und Linux:} Installiere die Version für dein System.
\end{itemize}

\subsection{Der Startbildschirm}
Auf einem Computer öffnet sich die App maximiert. Der Startbildschirm hat zwei Teile:

\begin{enumerate}
  \item \ui{Deine Bibliothek} (oben): die Bücher, die du heruntergeladen hast, das zuletzt
        gelesene zuerst, jeweils mit Fortschrittsbalken. Die Schaltfläche
        \ui{Bibliothek \& Regale} öffnet deine ganze Bibliothek.
  \item \ui{Kostenlose Bücher entdecken} (unten): ein Suchfeld, die Katalogsprache, die
        Kategorien, eine Reihe für jeden weiteren Katalog und zum Schluss der Katalog von Project
        Gutenberg.
\end{enumerate}

\shot{home}{Der Startbildschirm.}

Die beiden Schaltflächen oben rechts öffnen \ui{Anzeige \& Barrierefreiheit}
(Abschnitt~\ref{sec:display}) und \ui{Über}; dort steht, woher die Bücher kommen, und du findest
die Lizenzen der Schriften.

% ---------------------------------------------------------------------------
\section{Bücher finden}

\subsection{Suche}
Tippe in \ui{Nach Titel oder Autor suchen}. Die Suche
\begin{itemize}
  \item findet jeden Teil eines Titels oder Autorennamens: \emph{miserab} findet
        \emph{Les Misérables};
  \item ignoriert Akzente und Groß- und Kleinschreibung: \emph{misérable} und \emph{MISERABLE}
        liefern dasselbe;
  \item durchsucht alle Kataloge gleichzeitig; jeder zeigt seine Treffer in einer eigenen Reihe.
\end{itemize}
Leere das Feld mit der Schaltfläche \ui{×}, um wieder die beliebten Bücher zu sehen.

\subsection{Kategorien und Sprache}
Die Chips unter dem Suchfeld (\ui{Belletristik}, \ui{Krimi}, \ui{Science-Fiction}, \ui{Lyrik}…)
beschränken alle Kataloge auf ein Thema. Das Sprachmenü rechts zeigt nur Bücher in einer Sprache
oder \ui{Alle Sprachen}. Kataloge ohne Treffer für die gewählte Sprache oder Kategorie werden
ausgeblendet.

\subsection{Die Kataloge}
\begin{keytable}{Katalog}{Was er bietet}
  Project Gutenberg & Die größte kostenlose Bibliothek mit über 75.000 Büchern, nach
                      Beliebtheit sortiert. Viele gibt es als illustrierte Ausgabe. \\
  Standard Ebooks   & Sorgfältig korrigierte, schön gesetzte Klassiker auf Englisch. \\
  Wikisource        & Von Freiwilligen transkribierte Texte in vielen Sprachen. Das E-Book wird
                      beim Herunterladen erstellt; sehr lange Werke (etwa mehrbändige Romane)
                      sind eventuell nicht verfügbar. \\
  Internet Archive  & Scans alter gedruckter Ausgaben, die vor mindestens 95 Jahren erschienen
                      sind. Der Text wurde automatisch aus den Scans gewonnen und kann Fehler
                      enthalten. \\
\end{keytable}

\subsection{Die Seite eines Buches}
Tippe auf ein Cover, um die Seite des Buches zu öffnen. Sie zeigt die Zusammenfassung, Themen,
Sprache, Katalog und einen Hinweis zum Urheberrecht. Von dort kannst du:
\begin{itemize}
  \item \ui{Kostenlos herunterladen}: Das Buch wird auf deinem Gerät gespeichert und lässt sich
        offline lesen. Jedes Cover im Katalog hat außerdem eine kleine Download-Schaltfläche.
  \item \ui{Jetzt lesen} oder \ui{Weiterlesen}, sobald es heruntergeladen ist.
  \item \ui{Möchte ich lesen}, \ui{Gelesen} und \ui{Zu einem Regal hinzufügen}, auch schon vor
        dem Herunterladen (Abschnitt~\ref{sec:library}).
  \item \ui{Illustrierte Ausgabe laden}, für Gutenberg-Bücher, die du ohne Bilder geladen hast.
\end{itemize}

\shot{details}{Die Seite eines Buches.}

% ---------------------------------------------------------------------------
\section{Deine Bibliothek}\label{sec:library}

\ui{Bibliothek \& Regale} auf dem Startbildschirm öffnet alle Bücher, die du gespeichert hast:
heruntergeladene und solche, die du dir ohne Download für später gemerkt hast (mit blassem Cover
und einer kleinen Wolke).

\shot{library}{Die Bibliothek mit Ansichten und Regalen.}

\subsection{Ansichten, Filter und Sortierung}
Die Chips oben wählen eine Ansicht, jeweils mit der Anzahl der Bücher: \ui{Alle},
\ui{Heruntergeladen}, \ui{Möchte ich lesen}, \ui{Am Lesen}, \ui{Gelesen} und danach jedes deiner
Regale. Das Feld darüber filtert nach Titel oder Autor. Die Sortier-Schaltfläche (oben rechts)
ordnet die Bücher nach \ui{Zuletzt gelesen}, \ui{Titel}, \ui{Autor} (nach Nachname),
\ui{Fortschritt} oder \ui{Zuletzt hinzugefügt}. Die App merkt sich deine Wahl.

\subsection{Leselisten}
\begin{itemize}
  \item \ui{Möchte ich lesen} merkt ein Buch für später vor, ohne es herunterzuladen. Sobald du
        das Buch öffnest, verlässt es diese Liste.
  \item \ui{Gelesen} wird automatisch gesetzt, wenn du das Ende eines Buches erreichst (du kannst
        es rückgängig machen), oder von Hand.
\end{itemize}

\subsection{Regale}
Regale sind deine eigenen Sammlungen, etwa \emph{Lieblingsbücher} oder \emph{Philosophie}. Ein
Buch kann in mehreren Regalen stehen.
\begin{itemize}
  \item Lege eines mit \ui{Neues Regal} in der Bibliothek an, oder über
        \ui{Zu einem Regal hinzufügen} bei einem beliebigen Buch.
  \item Klicke mit der rechten Maustaste auf ein Regal (oder halte es gedrückt), um es
        umzubenennen oder zu löschen. Wenn du ein Regal löschst, bleiben seine Bücher in deiner
        Bibliothek.
\end{itemize}

\subsection{Das Buchmenü}
Klicke mit der rechten Maustaste auf ein Buch (oder halte es auf dem Telefon gedrückt). Du siehst
\ui{Buchdetails}, \ui{Zu einem Regal hinzufügen}, \ui{Möchte ich lesen},
\ui{Als gelesen markieren} und zwei Möglichkeiten, es zu entfernen:
\begin{itemize}
  \item \ui{Download entfernen} gibt Speicherplatz frei, lässt das Buch aber mit Notizen, Regalen
        und Leseposition in deiner Bibliothek.
  \item \ui{Aus der Bibliothek entfernen} löscht alles, auch Notizen und Lesezeichen.
\end{itemize}

% ---------------------------------------------------------------------------
\section{Lesen}

Tippe auf ein heruntergeladenes Buch, um es zu öffnen. Es öffnet sich dort, wo du aufgehört hast.

\shot{reader}{Der Reader mit seiner Werkzeugleiste.}

\subsection{Umblättern}
\begin{keytable}{Aktion}{So geht’s}
  Nächste Seite         & Rechten Rand antippen · \kright{} · \kbd{Bild ab} · \kbd{Leertaste} \\
  Vorherige Seite       & Linken Rand antippen · \kleft{} · \kbd{Bild auf} \\
  Ein Stück scrollen    & \kup{} \kdown{} oder das Mausrad \\
  Werkzeugleiste zeigen & Mitte der Seite antippen \\
\end{keytable}

\subsection{Die Werkzeugleiste}
Tippe auf die Mitte der Seite, um sie einzublenden. Oben: zurück zur Bibliothek, Suche
(Abschnitt~\ref{sec:search}), Lesezeichen, Inhalt und Anzeigeeinstellungen. Unten: vorheriges und
nächstes Kapitel, ein Schieberegler zum Springen zwischen Kapiteln und dein Fortschritt.

\subsection{Fortschritt und verbleibende Zeit}
Die unterste Zeile zeigt, wie viele Minuten das Kapitel noch dauert und wie weit du im Buch bist.
Die Zeit beruht auf deiner eigenen Lesegeschwindigkeit, die die App beim Lesen lernt.

\subsection{Inhalt und Bilder}
\ui{Inhalt} zeigt die echten Kapitel des Buches, auch wenn eine Datei mehrere davon enthält. Bilder
erscheinen im Text; tippe darauf, um sie im Vollbild zu sehen und zu vergrößern.

% ---------------------------------------------------------------------------
\section{Markierungen, Notizen und Lesezeichen}

\subsection{Text auswählen}
Auf dem Telefon hältst du ein Wort gedrückt und ziehst die Griffe. Am Computer ziehst du mit der
Maus über den Text oder doppelklickst auf ein Wort und klickst dann mit der \textbf{rechten
Maustaste} auf die Auswahl. Ein Menü erscheint mit:
\begin{itemize}
  \item \ui{Markieren}: hebt den Text in deiner Markierungsfarbe hervor.
  \item \ui{Notiz}: eine Markierung mit deiner eigenen Notiz, in einer von vier Farben.
  \item \ui{Nachschlagen} und \ui{Übersetzen} (Abschnitt~\ref{sec:lookup}).
  \item \ui{Kopieren}.
\end{itemize}
Markierungen rasten an ganzen Wörtern ein. Markierungen mit Notiz sind unterstrichen. Wähle eine
markierte Stelle erneut aus, um die Notiz zu bearbeiten oder die Markierung zu entfernen.

\subsection{Lesezeichen}
Nutze die Lesezeichen-Schaltfläche in der Werkzeugleiste oder \kbd{Strg+B}. Ein rotes Band zeigt,
dass die Seite ein Lesezeichen hat.

\subsection{Die Seitenleiste}
Die Inhalt-Schaltfläche öffnet eine Leiste mit drei Reitern: \ui{Inhalt}, \ui{Lesezeichen} und
\ui{Notizen}. Tippe auf einen Eintrag, um dorthin zu springen. \ui{Alles als Markdown kopieren}
kopiert alle Markierungen und Notizen, bereit zum Einfügen in ein anderes Programm.

% ---------------------------------------------------------------------------
\section{Wörterbuch und Übersetzung}\label{sec:lookup}

Wähle ein Wort (bis zu drei) aus und dann \ui{Nachschlagen}, um seine Bedeutungen zu sehen.
Bedeutungen in der Sprache des Buches stehen zuerst. Sie stammen aus dem Wiktionary und sind auf
Englisch verfasst.

Wähle eine beliebige Stelle aus und dann \ui{Übersetzen}. Wähle die Zielsprache aus der Liste; die
App merkt sie sich.

\begin{note}{Internet erforderlich}
Bedeutungen und Übersetzungen brauchen eine Internetverbindung. Der ausgewählte Text wird an das
Wiktionary oder an MyMemory gesendet; sonst nichts.
\end{note}

% ---------------------------------------------------------------------------
\section{In einem Buch suchen}\label{sec:search}

Drücke \kbd{Strg+F} oder die Such-Schaltfläche in der Werkzeugleiste und tippe mindestens zwei
Buchstaben. Die Suche ignoriert Groß- und Kleinschreibung und Akzente. Wähle einen Treffer, um
dorthin zu springen. Eine Leiste oben zeigt \emph{3 von 27}: Mit ihren Pfeilen oder mit \kbd{F3}
gehst du zum nächsten Treffer. \kbd{Esc} schließt sie.

% ---------------------------------------------------------------------------
\section{Anzeige und Barrierefreiheit}\label{sec:display}

Öffne die Einstellungen mit der Einstellungs-Schaltfläche auf dem Startbildschirm oder mit der
Text-Schaltfläche im Reader. Eine Vorschau oben zeigt jede Änderung sofort.

\shot{settings}{Anzeige \& Barrierefreiheit.}

\subsection{Sprache}
\ui{App-Sprache} stellt die Oberfläche sofort auf Deutsch, Englisch, Spanisch, Französisch,
Italienisch oder Portugiesisch um. \ui{Wie auf dem Gerät} folgt der Sprache deines Geräts.

\subsection{Text}
\begin{itemize}
  \item Design: \ui{Hell}, \ui{Sepia} oder \ui{Dunkel}.
  \item Schrift: vier Schriften, die in der App enthalten sind, damit Bücher auf jedem Gerät
        gleich aussehen:
  \begin{itemize}
    \item \textbf{Literata}, für das Lesen am Bildschirm entworfen (Standard);
    \item {\garamondfont EB Garamond}, eine klassische Buchschrift;
    \item {\sffamily Atkinson Hyperlegible}, so gestaltet, dass sich ähnliche Buchstaben leicht
          unterscheiden lassen;
    \item {\dyslexicfont OpenDyslexic}, für Leserinnen und Leser mit Legasthenie.
  \end{itemize}
  \item \ui{Schriftgröße}, \ui{Zeilenabstand} und \ui{Buchstabenabstand}.
  \item \ui{Blocksatz}: Schalte ihn aus, damit die Wortabstände gleich sind; viele Menschen mit
        Legasthenie finden das leichter.
\end{itemize}

\subsection{Seite}
\ui{Ränder} und \ui{Zeilenbreite}. Kürzere Zeilen lassen sich leichter verfolgen.

\subsection{Bildschirm}
\ui{Helligkeit} dunkelt die ganze App ab. \ui{Wärme} verringert das blaue Licht, wie ein
Nachtmodus. Beides wirkt zusätzlich zu den Einstellungen deines Geräts.

\ui{Auf Standard zurücksetzen} stellt alle diese Einstellungen wieder auf den Anfangszustand.

% ---------------------------------------------------------------------------
\section{Tastenkürzel}

\begin{keytable}{Tasten}{Aktion}
  \kright{} \kbd{Bild ab} \kbd{Leertaste} & Nächste Seite \\
  \kleft{} \kbd{Bild auf}                  & Vorherige Seite \\
  \kup{} \kdown{}                       & Ein Stück scrollen \\
  \kbd{Strg+F}                          & Im Buch suchen \\
  \kbd{F3}                              & Nächster Treffer \\
  \kbd{Strg+B}                          & Lesezeichen auf der Seite setzen oder entfernen \\
  \kbd{Esc}                             & Suche oder Leiste schließen, oder zurück zur
                                          Bibliothek \\
  \kbd{Umschalt} + Pfeiltasten          & Eine Textauswahl erweitern \\
\end{keytable}

% ---------------------------------------------------------------------------
\section{Datenschutz und deine Daten}

Illuminati hat keine Konten und kein Tracking. Deine Bibliothek, Notizen und Einstellungen werden
nur auf deinem Gerät gespeichert (unter Windows in
\texttt{\%APPDATA\%\textbackslash Illuminati}).

Die App verbindet sich nur mit dem Internet, um:
\begin{itemize}
  \item die Kataloge zu durchsuchen und Bücher und Cover herunterzuladen;
  \item ausgewählten Text nachzuschlagen oder zu übersetzen; er wird an das Wiktionary oder an
        MyMemory gesendet.
\end{itemize}

% ---------------------------------------------------------------------------
\section{Zum Urheberrecht}

Alle Bücher sind in den USA gemeinfrei. Einige Texte von Wikisource stehen stattdessen unter
freien Lizenzen (etwa CC BY-SA), die freies Lesen und Weitergeben erlauben. Aus dem Internet
Archive werden nur Bücher gezeigt, die vor mindestens 95 Jahren erschienen sind.

\begin{tip}{Prüfe die Gesetze deines Landes}
Das Urheberrecht ist von Land zu Land verschieden. Ein Buch, das in den USA frei ist, kann bei dir
noch geschützt sein. Prüfe deshalb vor dem Herunterladen die Gesetze deines Landes.
\end{tip}

% ---------------------------------------------------------------------------
\section{Fehlerbehebung}

\begin{keytable}{Meldung}{Was tun}
  Ein Katalog ist nicht erreichbar & Der Dienst ist vielleicht ausgelastet oder deine
                             Verbindung unterbrochen. Nutze später \ui{Erneut versuchen}; die
                             anderen Kataloge funktionieren weiter. \\
  Dieser Titel hat keine E-Book-Ausgabe & Wahrscheinlich ist es ein Hörbuch. Suche eine andere
                             Ausgabe. \\
  Wikisource konnte dieses Buch nicht erstellen & Das Werk ist zu lang für den Dienst von
                             Wikisource. Versuche einen seiner Bände oder einen anderen
                             Katalog. \\
  Tageslimit für kostenlose Übersetzungen erreicht & MyMemory erlaubt pro Tag nur eine begrenzte
                             Zahl kostenloser Übersetzungen. Versuch es morgen wieder. \\
  Der Download hat zu lange gedauert & Prüfe deine Verbindung und versuch es erneut. \\
\end{keytable}

% ---------------------------------------------------------------------------
\section{Danksagungen}

Bücher: Project Gutenberg, Standard Ebooks, Wikisource (über WS Export) und Internet Archive.
Bedeutungen: Wiktionary (CC BY-SA). Übersetzungen: MyMemory.

Schriften unter der SIL Open Font License: Literata, EB Garamond, Atkinson Hyperlegible (Braille
Institute) und OpenDyslexic (Abbie Gonzalez). Ihre Lizenzen findest du in der App unter
\ui{Über} → \ui{Lizenzen}.

\end{document}
